\documentclass[10pt,a4paper]{amsart}
\usepackage{amsmath,amssymb,amsthm,graphicx}
% Retained preamble for the Stacks Project chapter snapshot in this directory.
% The chapter's own preamble.tex is a copy of the shared Stacks Project preamble; the block
% below is sliced MECHANICALLY, line for line, from it. Nothing was added, removed or
% redefined; the wrapper class, the xr-hyper/hyperref wiring and the externaldocument list
% are omitted because the assembler supplies the document class.
%
% lines taken from preamble.tex: 3, 5, 9, 11, 12, 156, 157, 158, ... (51 lines total)
%
% package coverage restored from the chapter's own preamble (2026-10-13): xy, xy-2cell

\usepackage[all]{xy}\usepackage{verbatim}
\xyoption{2cell}
\UseAllTwocells

\newenvironment{reference}{\comment}{\endcomment}
\newenvironment{slogan}{\comment}{\endcomment}
\newenvironment{history}{\comment}{\endcomment}
\theoremstyle{plain}
\newtheorem{theorem}[subsection]{Theorem}
\newtheorem{proposition}[subsection]{Proposition}
\newtheorem{lemma}[subsection]{Lemma}
\theoremstyle{definition}
\newtheorem{definition}[subsection]{Definition}
\newtheorem{example}[subsection]{Example}
\newtheorem{exercise}[subsection]{Exercise}
\newtheorem{situation}[subsection]{Situation}
\theoremstyle{remark}
\newtheorem{remark}[subsection]{Remark}
\newtheorem{remarks}[subsection]{Remarks}
\numberwithin{equation}{subsection}
\def\lim{\mathop{\mathrm{lim}}\nolimits}
\def\colim{\mathop{\mathrm{colim}}\nolimits}
\def\Spec{\mathop{\mathrm{Spec}}}
\def\Hom{\mathop{\mathrm{Hom}}\nolimits}
\def\Ext{\mathop{\mathrm{Ext}}\nolimits}
\def\SheafHom{\mathop{\mathcal{H}\!\mathit{om}}\nolimits}
\def\SheafExt{\mathop{\mathcal{E}\!\mathit{xt}}\nolimits}
\def\Sch{\mathit{Sch}}
\def\Mor{\mathop{\mathrm{Mor}}\nolimits}
\def\Ob{\mathop{\mathrm{Ob}}\nolimits}
\def\Sh{\mathop{\mathit{Sh}}\nolimits}
\def\NL{\mathop{N\!L}\nolimits}
\def\CH{\mathop{\mathrm{CH}}\nolimits}
\def\proetale{{pro\text{-}\acute{e}tale}}
\def\etale{{\acute{e}tale}}
\def\QCoh{\mathit{QCoh}}
\def\Ker{\mathop{\mathrm{Ker}}}
\def\Im{\mathop{\mathrm{Im}}}
\def\Coker{\mathop{\mathrm{Coker}}}
\def\Coim{\mathop{\mathrm{Coim}}}
\DeclareMathSymbol{\boxtimes}{\mathbin}{AMSa}{"02}
\def\QCohstack{\mathcal{QC}\!\mathit{oh}}
\def\Cohstack{\mathcal{C}\!\mathit{oh}}
\def\Spacesstack{\mathcal{S}\!\mathit{paces}}
\def\Quotfunctor{\mathrm{Quot}}
\def\Hilbfunctor{\mathrm{Hilb}}
\def\Curvesstack{\mathcal{C}\!\mathit{urves}}
\def\Polarizedstack{\mathcal{P}\!\mathit{olarized}}
\def\Complexesstack{\mathcal{C}\!\mathit{omplexes}}
\def\Pic{\mathop{\mathrm{Pic}}\nolimits}
\def\Picardstack{\mathcal{P}\!\mathit{ic}}
\def\Picardfunctor{\mathrm{Pic}}
\def\Deformationcategory{\mathcal{D}\!\mathit{ef}}
\usepackage{cleveref}
\crefname{originallemma}{original lemma}{original lemmas}
\Crefname{originallemma}{Original lemma}{Original lemmas}
\makeatletter
\newlabel{lemma-cohomology-base-change}{{08IN}{source}{}{sourcepointer-1}{}}
\newlabel{lemma-cohomology-base-change@cref}{{[originallemma][0][]08IN}{[1][0][]source}{}{}{}}
\makeatother
\crefname{originallemma}{original lemma}{original lemmas}
\Crefname{originallemma}{Original lemma}{Original lemmas}
\makeatletter
\newlabel{lemma-compare-base-change}{{08IR}{source}{}{sourcepointer-2}{}}
\newlabel{lemma-compare-base-change@cref}{{[originallemma][0][]08IR}{[1][0][]source}{}{}{}}
\makeatother
\crefname{originallemma}{original lemma}{original lemmas}
\Crefname{originallemma}{Original lemma}{Original lemmas}
\makeatletter
\newlabel{lemma-countable}{{0CRV}{source}{}{sourcepointer-3}{}}
\newlabel{lemma-countable@cref}{{[originallemma][0][]0CRV}{[1][0][]source}{}{}{}}
\makeatother
\crefname{originallemma}{original lemma}{original lemmas}
\Crefname{originallemma}{Original lemma}{Original lemmas}
\makeatletter
\newlabel{lemma-countable-cohomology}{{0CRU}{source}{}{sourcepointer-4}{}}
\newlabel{lemma-countable-cohomology@cref}{{[originallemma][0][]0CRU}{[1][0][]source}{}{}{}}
\makeatother
\crefname{originallemma}{original lemma}{original lemmas}
\Crefname{originallemma}{Original lemma}{Original lemmas}
\makeatletter
\newlabel{lemma-quasi-coherence-pushforward-direct-sums}{{08FB}{source}{}{sourcepointer-5}{}}
\newlabel{lemma-quasi-coherence-pushforward-direct-sums@cref}{{[originallemma][0][]08FB}{[1][0][]source}{}{}{}}
\makeatother
\crefname{originallemma}{original lemma}{original lemmas}
\Crefname{originallemma}{Original lemma}{Original lemmas}
\makeatletter
\newlabel{sites-cohomology-lemma-dual-perfect-complex}{{08JJ}{source}{}{sourcepointer-6}{}}
\newlabel{sites-cohomology-lemma-dual-perfect-complex@cref}{{[originallemma][0][]08JJ}{[1][0][]source}{}{}{}}
\makeatother
\crefname{originallemma}{original lemma}{original lemmas}
\Crefname{originallemma}{Original lemma}{Original lemmas}
\makeatletter
\newlabel{spaces-limits-lemma-descend-finite-presentation}{{07SK}{source}{}{sourcepointer-7}{}}
\newlabel{spaces-limits-lemma-descend-finite-presentation@cref}{{[originallemma][0][]07SK}{[1][0][]source}{}{}{}}
\makeatother
\crefname{originallemma}{original lemma}{original lemmas}
\Crefname{originallemma}{Original lemma}{Original lemmas}
\makeatletter
\newlabel{spaces-limits-lemma-descend-flat}{{08K0}{source}{}{sourcepointer-8}{}}
\newlabel{spaces-limits-lemma-descend-flat@cref}{{[originallemma][0][]08K0}{[1][0][]source}{}{}{}}
\makeatother
\title{Stacks, Derived Categories of Spaces: computing-sections-as-colim}
\author{The Stacks Project Authors}
\date{stacks-project:spaces-perfect.tex:local-snapshot-20261012; GNU FDL 1.2 or later}
\begin{document}
\maketitle
\noindent\textit{Editorial note.} Statement span is lines 7022-7040 of sources/stacks-spaces-perfect/pinned-original-source.tex; proof span is lines 7042-7085, the complete author proof (42 body lines).\par\medskip
\noindent\textit{Editorial note.} Navigation stubs: 8 cited-but-not-included labels are supplied as non-load-bearing pointers with their pinned tags where the snapshot tag index has them.\par\medskip
\noindent\textit{Editorial note.} extra\_packages is empty on purpose: the retained preamble already loads xy with options and a second declaration clashes in the assembler wrapper.\par\medskip
\noindent\textit{Editorial note.} Text is retained verbatim, including the author's own typos and placeholders.\par\medskip
\noindent\textit{Editorial note.} Queue index 49 of queue\_spaceperfect.json (pb=2028, stmt 7022-7040, proof 7042-7085).\par\medskip
\noindent\textit{Editorial note.} stacks-spaces-perfect wave sp01 (2026-10-14): fresh adjudication accept per lane-s3/sp01\_judgments; pool per Chief 12:4x unfreeze (reserve-pool workflow, six-check clean); spaces-perfect.tex 126 with-proof rows, 70 b107-occupied labels avoided, 56 free rows judged 32 accept / 24 reject.\par\medskip
\noindent\textit{Editorial note.} Cross-chapter same-label warning (NOT an occupancy conflict, Chief ruling 2026-10-13 evening): label 'lemma-computing-sections-as-colim' also exists in stacks-perfect — a different proposition in a different chapter file.\par\medskip
\noindent\textit{Allowed standard prerequisites (editorial).}\begin{itemize}
\item Explicit prerequisite boundary. The cited inputs are prerequisites; the argument in the retained span is not replaced by any cited result.
\end{itemize}\medskip
\ifdefined\hypertarget\hypertarget{sourcepointer-1}{}\fi
\noindent\textit{Source reference.} Original lemma 08IN: sources/stacks-spaces-perfect/pinned-original-source.tex, source lines 4883--4892. Cited by the retained statement or proof but not part of the collected spans; supplied only as a navigation pointer. The argument does not depend on its content.\par\medskip
\ifdefined\hypertarget\hypertarget{sourcepointer-2}{}\fi
\noindent\textit{Source reference.} Original lemma 08IR: sources/stacks-spaces-perfect/pinned-original-source.tex, source lines 5080--5095. Cited by the retained statement or proof but not part of the collected spans; supplied only as a navigation pointer. The argument does not depend on its content.\par\medskip
\ifdefined\hypertarget\hypertarget{sourcepointer-3}{}\fi
\noindent\textit{Source reference.} Original lemma 0CRV: sources/stacks-spaces-perfect/pinned-original-source.tex, source lines 6969--6980. Cited by the retained statement or proof but not part of the collected spans; supplied only as a navigation pointer. The argument does not depend on its content.\par\medskip
\ifdefined\hypertarget\hypertarget{sourcepointer-4}{}\fi
\noindent\textit{Source reference.} Original lemma 0CRU: sources/stacks-spaces-perfect/pinned-original-source.tex, source lines 6932--6944. Cited by the retained statement or proof but not part of the collected spans; supplied only as a navigation pointer. The argument does not depend on its content.\par\medskip
\ifdefined\hypertarget\hypertarget{sourcepointer-5}{}\fi
\noindent\textit{Source reference.} Original lemma 08FB: sources/stacks-spaces-perfect/pinned-original-source.tex, source lines 863--868. Cited by the retained statement or proof but not part of the collected spans; supplied only as a navigation pointer. The argument does not depend on its content.\par\medskip
\ifdefined\hypertarget\hypertarget{sourcepointer-6}{}\fi
\noindent\textit{Source reference.} Original lemma 08JJ: sources/stacks-spaces-perfect/pinned-original-source.tex, source lines 7085--7085. Cited by the retained statement or proof but not part of the collected spans; supplied only as a navigation pointer. The argument does not depend on its content.\par\medskip
\ifdefined\hypertarget\hypertarget{sourcepointer-7}{}\fi
\noindent\textit{Source reference.} Original lemma 07SK: sources/stacks-spaces-perfect/pinned-original-source.tex, source lines 7085--7085. Cited by the retained statement or proof but not part of the collected spans; supplied only as a navigation pointer. The argument does not depend on its content.\par\medskip
\ifdefined\hypertarget\hypertarget{sourcepointer-8}{}\fi
\noindent\textit{Source reference.} Original lemma 08K0: sources/stacks-spaces-perfect/pinned-original-source.tex, source lines 7085--7085. Cited by the retained statement or proof but not part of the collected spans; supplied only as a navigation pointer. The argument does not depend on its content.\par\medskip
\begin{lemma}
\label{lemma-computing-sections-as-colim}
Let $A$ be a ring. Let $f : U \to X$ be a flat morphism of algebraic spaces
of finite presentation over $A$. Then
\begin{enumerate}
\item there exists an inverse system of perfect objects $L_n$ of
$D(\mathcal{O}_X)$ such that
$$
R\Gamma(U, Lf^*K) = \text{hocolim}\ R\Hom_X(L_n, K)
$$
in $D(A)$ functorially in $K$ in $D_\QCoh(\mathcal{O}_X)$, and
\item there exists a system of perfect objects $E_n$ of
$D(\mathcal{O}_X)$ such that
$$
R\Gamma(U, Lf^*K) = \text{hocolim}\ R\Gamma(X, E_n \otimes^\mathbf{L} K)
$$
in $D(A)$ functorially in $K$ in $D_\QCoh(\mathcal{O}_X)$.
\end{enumerate}
\end{lemma}

\begin{proof}
By Lemma \ref{lemma-cohomology-base-change} we have
$$
R\Gamma(U, Lf^*K) = R\Gamma(X, Rf_*\mathcal{O}_U \otimes^\mathbf{L} K)
$$
functorially in $K$. Observe that $R\Gamma(X, -)$ commutes with
homotopy colimits because it commutes with direct sums by
Lemma \ref{lemma-quasi-coherence-pushforward-direct-sums}.
Similarly, $- \otimes^\mathbf{L} K$ commutes with derived colimits
because $- \otimes^\mathbf{L} K$ commutes with direct sums
(because direct sums in $D(\mathcal{O}_X)$
are given by direct sums of representing complexes).
Hence to prove (2) it suffices to write
$Rf_*\mathcal{O}_U = \text{hocolim} E_n$ for a system of
perfect objects $E_n$ of $D(\mathcal{O}_X)$. Once this is done
we obtain (1) by setting $L_n = E_n^\vee$, see Cohomology on Sites,
Lemma \ref{sites-cohomology-lemma-dual-perfect-complex}.

\medskip\noindent
Write $A = \colim A_i$ with $A_i$ of finite type over $\mathbf{Z}$. By
Limits of Spaces, Lemma \ref{spaces-limits-lemma-descend-finite-presentation}
we can find an $i$ and morphisms $U_i \to X_i \to \Spec(A_i)$
of finite presentation whose base change to $\Spec(A)$ recovers
$U \to X \to \Spec(A)$.
After increasing $i$ we may assume that $f_i : U_i \to X_i$ is
flat, see Limits of Spaces, Lemma
\ref{spaces-limits-lemma-descend-flat}.
By Lemma \ref{lemma-compare-base-change}
the derived pullback of $Rf_{i, *}\mathcal{O}_{U_i}$
by $g : X \to X_i$ is equal to $Rf_*\mathcal{O}_U$.
Since $Lg^*$ commutes with derived colimits, it suffices
to prove what we want for $f_i$. Hence we may assume that
$U$ and $X$ are of finite type over $\mathbf{Z}$.

\medskip\noindent
Assume $f : U \to X$ is a morphism of algebraic spaces
of finite type over $\mathbf{Z}$. To finish the proof
we will show that $Rf_*\mathcal{O}_U$ is a homotopy
colimit of perfect complexes. To see this we apply Lemma \ref{lemma-countable}.
Thus it suffices to show that $R^if_*\mathcal{O}_U$
has countable sets of sections over affines \'etale over $X$.
This follows from Lemma \ref{lemma-countable-cohomology}
applied to the structure sheaf.
\end{proof}

\end{document}
